% Source: https://artofproblemsolving.com/community/c3h1290706_formula_marathon
% Stable identity: aops-topic-1290706-q01
A picture explains everything: take a $30$-$60$-$90$ triangle $ABC$ with $AC=1$, $BC=\sqrt3$, $AB=2$, and extend $CB$ to $D$ with $BD=BA=2$. Triangle $ABD$ is isosceles with exterior angle $30^\circ$ at $B$, so $\angle ADC=15^\circ$.

\begin{asy}
usepackage("amsmath");
size(6.5cm);
defaultpen(fontsize(10pt));
path ra(pair P,pair u,pair v,real s=0.2){return P+s*unit(u)--P+s*unit(u)+s*unit(v)--P+s*unit(v);}
real s3=sqrt(3);
pair C=(0,0), A=(0,1), B=(-s3,0), D=(-2-s3,0);
draw(D--C--A--cycle,linewidth(0.9)); draw(A--B,blue);
draw(ra(C,D-C,A-C,0.18));
label("$C$",C,SE); label("$A$",A,N); label("$B$",B,S); label("$D$",D,SW);
label("$1$",(A+C)/2,E); label("$\sqrt{3}$",(B+C)/2,S); label("$2$",(B+D)/2,S); label("$2$",B+0.72*(A-B),SE,blue);
draw(arc(D,1.1,0,15)); label("$15^\circ$",D+1.6*dir(7.5),fontsize(8pt));
draw(arc(B,0.4,0,30),blue); label("$30^\circ$",B+0.8*dir(15),blue+fontsize(8pt));
label("\mbox{}",(min(currentpicture,user=true)+max(currentpicture,user=true))/2,black); // keeps the text after the figure black when LaTeX typesets the labels inline
\end{asy}

In the right triangle $ACD$: $AC=1$, $DC=2+\sqrt3$, so $\tan15^\circ=\dfrac{1}{2+\sqrt3}=2-\sqrt3$, and
\[
AD^2=1+\left(2+\sqrt3\right)^2=8+4\sqrt3=\left(\sqrt6+\sqrt2\right)^2,\qquad AD=\sqrt6+\sqrt2 .
\]

(a) $\sin15^\circ=\dfrac{AC}{AD}=\dfrac{1}{\sqrt6+\sqrt2}=\boxed{\dfrac{\sqrt6-\sqrt2}{4}}$. This agrees with the given form: from $\tan15^\circ=2-\sqrt3$ and $1+\cot^2\theta=\csc^2\theta$ one gets $\sin^2 15^\circ=\dfrac{1}{1+\left(2+\sqrt3\right)^2}=\dfrac{1}{8+4\sqrt3}=\dfrac{2-\sqrt3}{4}$, so $\sin15^\circ=\dfrac{\sqrt{2-\sqrt3}}{2}$; and $\left(\dfrac{\sqrt6-\sqrt2}{2}\right)^2=\dfrac{8-2\sqrt{12}}{4}=2-\sqrt3$, so $\sqrt{2-\sqrt3}=\dfrac{\sqrt6-\sqrt2}{2}$.

(b) $\cos15^\circ=\dfrac{DC}{AD}=\dfrac{2+\sqrt3}{\sqrt6+\sqrt2}=\left(2+\sqrt3\right)\cdot\dfrac{\sqrt6-\sqrt2}{4}=\dfrac{2\sqrt6-2\sqrt2+3\sqrt2-\sqrt6}{4}=\boxed{\dfrac{\sqrt6+\sqrt2}{4}}$.

(c) Taking reciprocals, and using $\sin75^\circ=\cos15^\circ$, $\cos75^\circ=\sin15^\circ$:
\[
\csc15^\circ=\sec75^\circ=\dfrac{4}{\sqrt6-\sqrt2}=\boxed{\sqrt6+\sqrt2},\qquad
\sec15^\circ=\csc75^\circ=\dfrac{4}{\sqrt6+\sqrt2}=\boxed{\sqrt6-\sqrt2},
\]
\[
\cot15^\circ=\tan75^\circ=\boxed{2+\sqrt3},\qquad \cot75^\circ=\tan15^\circ=\boxed{2-\sqrt3}.
\]
